% Cours autonome : convention LaTeX du 8 octobre 2026 (12 pt, interligne 1,5,
% chapitres sur une nouvelle page), sans le préambule 11 pt des anciens cours.
\documentclass[12pt,a4paper,oneside]{report}
\usepackage[T1]{fontenc}
\usepackage[utf8]{inputenc}
\usepackage[french]{babel}
\usepackage{lmodern}
\usepackage{microtype}
\usepackage[a4paper,margin=2.35cm]{geometry}
\usepackage{setspace}
\onehalfspacing
\setlength{\parindent}{0pt}
\setlength{\parskip}{0.5\baselineskip}
\usepackage{amsmath,amssymb,mathtools}
\usepackage{booktabs,longtable,array}
\usepackage{xcolor}
\usepackage{hyperref}
\usepackage{tikz}
\usetikzlibrary{arrows.meta,positioning,calc}
\hypersetup{colorlinks=true,linkcolor=blue!55!black,citecolor=blue!55!black,urlcolor=blue!55!black,
pdftitle={Cours 12 -- Vibrance et saturation : mesurer et reproduire Lightroom},
pdfauthor={Projet IPR -- Scholar}}
\definecolor{iprblue}{HTML}{1F4E79}
\definecolor{iprgreen}{HTML}{2D6A4F}
\newcommand{\IPR}{\textsc{ipr}}
\newcommand{\resultat}[1]{\paragraph{À retenir.} #1}
\clubpenalty=10000
\widowpenalty=10000
\raggedbottom

\title{Cours 12 -- Vibrance et saturation\\
\large Mesurer Lightroom pixel par pixel et en faire un modèle}
\author{Projet \IPR{} -- documentation Scholar}
\date{22 octobre 2026\\[3em]
\small Du curseur à l'équation : espaces de couleur, Oklab, mesure d'un logiciel de
référence, tables de réponse, rotation de teinte, validation et implémentation.\\
Valeurs issues des campagnes du 21 et 22 octobre 2026 sur 49 scènes.}

\begin{document}
\maketitle
\tableofcontents

\chapter*{Comment utiliser ce cours}
\addcontentsline{toc}{chapter}{Comment utiliser ce cours}
Ce cours explique comment \IPR{} reproduit les curseurs Saturation et
Vibrance de Lightroom. Il ne se contente pas de donner les formules : il
montre la démarche qui y a conduit. On mesure d'abord ce que fait le logiciel
de référence, on vérifie quelles grandeurs expliquent son comportement, puis
on construit un modèle et on le valide sur des images qui n'ont pas servi à
le construire.

Les chapitres 1 et 2 posent le problème et les outils colorimétriques. Le
chapitre 3 décrit la mesure de Lightroom et ce qu'elle révèle. Les
chapitres 4 et 5 construisent le modèle, gain puis rotation de teinte. Le
chapitre 6 le valide, le chapitre 7 décrit son implémentation, ses limites,
et propose des exercices corrigés.

\paragraph{Prérequis.} Calcul matriciel $3\times3$, logarithme, notions de
statistique (médiane, variance). Les espaces de couleur sont rappelés au
chapitre 2.

\paragraph{Périmètre.} Les nombres cités proviennent des preuves JSON de
\texttt{docs/validation/raw} (campagnes \texttt{mesure-couleur} et
\texttt{modele-vibrance-lr}) ; elles font autorité en cas d'écart avec ce
document. Les tables décrivent Lightroom avec le profil Camera Neutral, sur les
boîtiers du corpus (Canon R6 Mark III, R7 et EOS 60D) ; elles ne constituent
pas une vérité universelle sur Lightroom.

\chapter{Saturation et vibrance : ce que l'on cherche à reproduire}

\section{Deux curseurs voisins, deux intentions}
Les deux curseurs modifient l'intensité des couleurs, sans changer la
luminosité générale de l'image. Leur intention diffère pourtant.

\begin{itemize}
\item La \textbf{saturation} agit uniformément : chaque couleur voit son
intensité multipliée par le même facteur. À $-100$, l'image devient grise ; à
$+100$, toutes les couleurs, y compris celles qui sont déjà vives, sont
renforcées autant.
\item La \textbf{vibrance} est sélective. Elle renforce surtout les couleurs
ternes et épargne celles qui sont déjà intenses. Les tons chair sont réputés
protégés. Le photographe s'en sert pour réveiller une image sans faire
« baver » les couleurs vives.
\end{itemize}

La description de la vibrance reste qualitative : les éditeurs ne publient pas
leur formule. Pour reproduire Lightroom, il faut donc le \emph{mesurer}.

\section{Pourquoi un modèle mathématique simple ne suffit pas}
Un premier modèle de vibrance, utilisé par \IPR{} jusqu'en octobre 2026,
pondérait le gain par la saturation relative du pixel,
$g_v=1+(V/100)(1-c)$ avec $c=(\max-\min)/\max$. Un second prototype (A1/A2)
travaillait en Oklab avec un poids $(1-q)^2$ et une protection des tons chair.
Mesurés sur 49 scènes, ces modèles restaient loin de Lightroom :

\begin{itemize}
\item à $+25$ et $+50$, ils renforçaient environ deux fois trop les couleurs
peu saturées ;
\item à $+100$, ils renforçaient trop les couleurs saturées ;
\item à $-100$, ils retiraient trop d'intensité aux couleurs saturées ;
\item ils ne modifiaient jamais la clarté, alors que Lightroom assombrit
légèrement en désaturant, ni la teinte, alors que Lightroom la fait tourner.
\end{itemize}

Le dernier point est le plus parlant : aucun choix de constantes n'aurait
corrigé un modèle qui ne possède pas le degré de liberté nécessaire.

\resultat{On ne devine pas le comportement d'un logiciel de référence : on le
mesure, on identifie les variables qui l'expliquent, puis on choisit la
représentation du modèle.}

\section{La démarche suivie}
\begin{enumerate}
\item \textbf{Constituer des paires} : pour chaque photographie, un export
Lightroom neutre et un export avec le seul curseur modifié.
\item \textbf{Comparer pixel par pixel} chaque export à son neutre : gain
d'intensité, changement de clarté, rotation de teinte.
\item \textbf{Chercher les variables explicatives} : de quoi dépend la réponse
d'un pixel ? De sa saturation, de sa clarté, de sa teinte, de l'image ?
\item \textbf{Construire un modèle} à partir de ces variables.
\item \textbf{Valider} sur des images exclues de la construction, puis sur le
rendu réel d'\IPR{}.
\end{enumerate}

Ces cinq étapes structurent les chapitres suivants.

\chapter{Les outils colorimétriques}

\section{RGB linéaire et changements de primaires}
\IPR{} travaille sur un master en RGB linéaire Rec.2020, D65 : les valeurs sont
proportionnelles à la lumière, sans courbe de transfert. Passer d'un jeu de
primaires à un autre est alors une simple application linéaire,
\begin{equation}
\mathbf{p}_{\text{sRGB}} = M\,\mathbf{p}_{\text{Rec.2020}},\qquad
M = M_{\text{sRGB}\to XYZ}^{-1}\,M_{\text{Rec.2020}\to XYZ},
\end{equation}
où chaque matrice $M_{\cdot\to XYZ}$ se déduit des coordonnées $xy$ des
primaires et du blanc~\cite{itu2020,iec61966}. Le passage n'écrête rien : une
couleur Rec.2020 hors du gamut sRGB obtient simplement une composante négative,
et l'aller-retour $M^{-1}M$ restitue la valeur exacte.

\resultat{Les tables de vibrance ont été mesurées sur des exports sRGB : \IPR{}
convertit le master en sRGB linéaire, applique le modèle, puis revient en
Rec.2020, sans écrêtage dans l'intervalle.}

\section{Pourquoi Oklab}
Pour agir sur « l'intensité d'une couleur » sans toucher à sa luminosité ni à
sa teinte, il faut un espace où ces trois notions sont séparées. Oklab,
proposé par Björn Ottosson en 2020~\cite{ottosson2020}, est construit pour que
\begin{itemize}
\item $L$ représente la clarté perçue ;
\item le plan $(a,b)$ porte la couleur : la distance au gris est la
\textbf{chroma} $C=\sqrt{a^2+b^2}$ et l'angle est la \textbf{teinte}
$h=\operatorname{atan2}(b,a)$ ;
\item une variation de $C$ à $L$ et $h$ fixes soit perçue comme une variation
d'intensité sans changement de teinte, mieux que dans HSV ou CIELAB.
\end{itemize}

Depuis le RGB linéaire sRGB, la transformation compte deux matrices et une
racine cubique :
\begin{align}
(l,m,s)^\top &= M_1\,(R,G,B)^\top,\\
(L,a,b)^\top &= M_2\,(\sqrt[3]{l},\sqrt[3]{m},\sqrt[3]{s})^\top,
\end{align}
avec
\begin{equation}
M_1=\begin{pmatrix}
0{,}4122 & 0{,}5363 & 0{,}0514\\
0{,}2119 & 0{,}6807 & 0{,}1074\\
0{,}0883 & 0{,}2817 & 0{,}6300\end{pmatrix},\quad
M_2=\begin{pmatrix}
0{,}2105 & 0{,}7936 & -0{,}0041\\
1{,}9780 & -2{,}4286 & 0{,}4506\\
0{,}0259 & 0{,}7828 & -0{,}8087\end{pmatrix}.
\end{equation}
\IPR{} utilise la racine cubique signée ($\sqrt[3]{-x}=-\sqrt[3]{x}$) : une
composante négative ou supérieure à 1, fréquente dans un master HDR, reste
calculable. Le retour inverse les deux étapes.

\section{Deux mesures d'intensité : chroma et saturation}
La \textbf{chroma} Oklab $C$ est une mesure absolue : elle grandit avec la
distance au gris. Une couleur sombre et une couleur claire de même aspect
« vif » n'ont pas la même chroma.

La \textbf{saturation} est une mesure relative. Le modèle utilise celle d'un
espace RGB à très large gamut, ProPhoto RGB (ROMM, ISO 22028-2~\cite{iso22028}) :
\begin{equation}
s=\frac{\max(R,G,B)-\min(R,G,B)}{\max(R,G,B)}\in[0,1],
\end{equation}
calculée en ProPhoto \emph{linéaire}. Le passage sRGB vers ProPhoto change de
blanc (D65 vers D50) par la transformation de Bradford~\cite{lam1985}.
$s=0$ désigne un gris, $s=1$ une couleur sur le bord du gamut ProPhoto.

Pourquoi ProPhoto ? Parce que c'est, parmi les variables testées, celle qui
explique le mieux le comportement positif de Lightroom (chapitre 3), ce qui
est cohérent avec un espace de travail interne fondé sur les primaires
ProPhoto. Ce n'est pas une affirmation sur le code de Lightroom, mais un
constat de mesure.

\resultat{Trois coordonnées décrivent un pixel pour la vibrance : sa
saturation ProPhoto $s$, sa clarté Oklab $L$ et sa teinte Oklab $h$. La chroma
$C$ sert à mesurer l'effet et à repérer les quasi-gris.}

\section{Rotation dans le plan de couleur}
Faire tourner la teinte d'un angle $\theta$ sans changer la chroma revient à
appliquer la matrice de rotation plane au couple $(a,b)$ :
\begin{equation}
\begin{pmatrix}a'\\ b'\end{pmatrix}=
\begin{pmatrix}\cos\theta & -\sin\theta\\ \sin\theta & \cos\theta\end{pmatrix}
\begin{pmatrix}a\\ b\end{pmatrix}.
\end{equation}
Cette matrice est orthogonale : elle conserve la norme, donc $C$, et ajoute
$\theta$ à $h$. Combinée à un gain $g$ sur la chroma et à un gain $k$ sur la
clarté, la transformation complète d'un pixel s'écrit
\begin{equation}
L'=k\,L,\qquad (a',b')^\top = g\,R(\theta)\,(a,b)^\top .
\label{eq:transformation}
\end{equation}
Gain et rotation commutent ($gR=Rg$) et restent indépendants : l'un règle
l'intensité, l'autre la teinte.

\resultat{Ce n'est pas une matrice unique appliquée à toute l'image : $g$, $k$
et $\theta$ dépendent du pixel et de la valeur du curseur. La transformation
globale n'est donc pas linéaire, même si elle l'est localement dans le plan
$(a,b)$.}

\chapter{Mesurer Lightroom}

\section{Le protocole}
\subsection{Des paires d'exports comparables}
Chaque scène est exportée de Lightroom en TIFF 16 bits sRGB, une fois à l'état
neutre et une fois par valeur testée du curseur : $-100$, $-50$, $-25$, $+25$,
$+50$, $+100$. Les préréglages utilisés sont \emph{autonomes} : chacun fixe
tout l'état de développement (profil Camera Neutral, balance des blancs Prise
de vue, corrections optiques, bruit et netteté désactivés, tous les autres
curseurs à zéro) et ne change que le curseur étudié. Appliqué à n'importe
quelle photo, un préréglage donne donc un état connu, quels que soient ses
réglages précédents.

L'analyseur ne se fie pas au nom des fichiers. Il lit le XMP embarqué dans
chaque TIFF, y trouve la valeur du curseur, et associe l'export au neutre dont
\emph{tous les autres réglages} sont identiques. Une association ambiguë, un
état manquant ou dupliqué arrête l'analyse.

\subsection{Comparer un pixel à lui-même}
Dans une paire, le neutre et la variante ont la même géométrie : le pixel
$(x,y)$ de l'un correspond exactement au pixel $(x,y)$ de l'autre. On peut
donc mesurer, pour chaque pixel, l'effet du curseur :
\begin{equation}
\gamma = \ln\frac{C_{\text{variante}}}{C_{\text{neutre}}},\qquad
\lambda = \ln\frac{L_{\text{variante}}}{L_{\text{neutre}}},\qquad
\Delta h = h_{\text{variante}}-h_{\text{neutre}}\ (\text{ramené à }]-180^\circ,180^\circ]).
\end{equation}
Le logarithme rend les gains additifs et symétriques : doubler la chroma donne
$\gamma=\ln 2\approx0{,}69$, la diviser par deux donne $-0{,}69$.

\subsection{Échantillonnage et exclusions}
Un pixel sur 8 dans chaque direction est lu, puis 40\,000 à 60\,000 pixels par
état sont tirés au hasard avec une graine fixe (reproductible). Sont exclus :
\begin{itemize}
\item les pixels dont un canal est écrêté ($\le0{,}002$ ou $\ge0{,}995$ en
valeur encodée) dans l'un des deux exports : leur valeur ne dit plus rien de
l'effet du curseur ;
\item les pixels très sombres ($L\le0{,}08$), dominés par le bruit ;
\item pour les tables, les quasi-gris ($C\le0{,}01$), traités à part
(section~\ref{sec:quasigris}).
\end{itemize}

\resultat{Comparer chaque export à son propre neutre, et non à \IPR{},
mesure le comportement du curseur indépendamment du rendu de base.}

\section{Quelle variable explique la réponse ?}
On range les pixels en 20 classes d'une variable candidate, on retire de
chaque pixel la médiane de sa classe, et on mesure la variance restante de
$\gamma$. Plus elle est basse, mieux la variable explique la réponse.

\begin{center}
\begin{tabular}{lrrrr}
\toprule
Valeur & variance totale & $C/L$ Oklab & $C$ Oklab & $s$ ProPhoto\\
\midrule
$-100$ & 0,7465 & 0,5841 & \textbf{0,2031} & 0,5246\\
$-50$ & 0,0115 & 0,0078 & \textbf{0,0037} & 0,0054\\
$+50$ & 0,0184 & 0,0130 & 0,0183 & \textbf{0,0130}\\
$+100$ & 0,1355 & \textbf{0,0715} & 0,1344 & 0,0759\\
\bottomrule
\end{tabular}
\end{center}

Au positif, la saturation relative ($s$ ou $C/L$) explique bien plus que la
chroma absolue. Au négatif, la chroma absolue explique mieux, mais aucune
variable seule ne suffit.

On ajoute ensuite des variables, en croisant les classes. La part de variance
expliquée devient :

\begin{center}
\begin{tabular}{lrrrrr}
\toprule
Valeur & $s$ & $s$+$L$ & $s$+$L$+$h$ & $s$+scène & $s$+$L$+$h$+scène\\
\midrule
$-100$ & 73\,\% & 79\,\% & 85\,\% & 80\,\% & 88\,\%\\
$+25$ & 25\,\% & 52\,\% & 93\,\% & 56\,\% & 95\,\%\\
$+50$ & 29\,\% & 59\,\% & 94\,\% & 58\,\% & 96\,\%\\
$+100$ & 44\,\% & 77\,\% & 95\,\% & 67\,\% & 97\,\%\\
\bottomrule
\end{tabular}
\end{center}

Saturation, clarté et teinte expliquent 93 à 95\,\% de la variance au positif.
Connaître en plus l'image d'où vient le pixel n'apporte que 2\,\%.

\resultat{La vibrance de Lightroom est une fonction fixe du pixel, pas une
adaptation à l'image. Les grandes différences constatées d'une scène à
l'autre viennent de leur contenu : une scène riche en tons moyens peu saturés
réagit davantage.}

\section{La forme de la réponse}
\subsection{Le gain selon la saturation et la clarté}
Gain médian de chroma à $+100$, selon la saturation $s$ (lignes) et la clarté
$L$ (colonnes) du pixel neutre :

\begin{center}
\begin{tabular}{lrrrrrr}
\toprule
$s$ \textbackslash{} $L$ & 0,19 & 0,38 & 0,53 & 0,68 & 0,82 & 0,96\\
\midrule
0,00--0,10 & -- & 2,80 & 3,73 & 3,07 & 2,06 & 1,72\\
0,10--0,20 & 1,53 & 2,89 & 3,70 & 3,11 & 2,13 & 1,82\\
0,20--0,30 & 1,37 & 2,55 & 3,35 & 2,88 & 2,20 & 1,89\\
0,30--0,45 & 1,24 & 2,00 & 2,42 & 2,44 & 2,13 & 1,85\\
0,45--0,60 & 1,22 & 1,67 & 1,74 & 1,83 & 1,92 & 1,79\\
0,60--0,80 & 1,19 & 1,38 & 1,42 & 1,45 & 1,53 & --\\
0,80--1,00 & 1,13 & 1,22 & 1,31 & -- & -- & --\\
\bottomrule
\end{tabular}
\end{center}

Trois tendances se lisent :
\begin{itemize}
\item le gain décroît quand la saturation augmente : les couleurs déjà vives
sont protégées, c'est la définition même de la vibrance ;
\item il culmine sur les tons moyens ($L\approx0{,}5$) et retombe fortement
dans les ombres et les hautes lumières ;
\item la dépendance à $L$ est aussi forte que celle à $s$. L'ancien modèle, qui
l'ignorait, ne pouvait pas reproduire Lightroom.
\end{itemize}

\subsection{Une réponse progressive au curseur}
Au maximum du gain, la chroma est multipliée par 1,18 à $+25$, 1,42 à $+50$ et
2,49 à $+100$. En logarithme, $0{,}16$, $0{,}35$ et $0{,}91$ : la réponse
croît plus vite que le curseur, à peu près comme $V^{1{,}3}$. Le curseur est
doux au début et s'accélère vers $+100$.

\subsection{Le côté négatif}
À $-25$ et $-50$, le gain est presque uniforme (0,70 puis 0,45) jusqu'à
$s\approx0{,}5$ ; au-delà, les couleurs saturées sont protégées. À $-100$, les
couleurs faibles sont presque entièrement désaturées (gain 0,03 à 0,05), les
plus saturées gardent jusqu'à 0,35 de leur chroma.

\subsection{La clarté change aussi}
Lightroom ne garde pas la clarté constante. Pour les couleurs saturées, à
$-100$, $L$ est multipliée par 0,88 en médiane : désaturer assombrit. Ces
gains de clarté sont mesurés et reproduits comme ceux de la chroma. Ils
dépendent eux aussi de la couleur : certains bleus s'éclaircissent légèrement
en désaturant.

\subsection{Les quasi-gris}
\label{sec:quasigris}
Les pixels presque gris ($C<0{,}01$), souvent du bruit chromatique, reçoivent un
traitement particulier. À $+100$, Lightroom ne multiplie leur chroma que par
1,2 environ, quand leur classe de saturation prévoirait 1,55. À $-100$, au
contraire, il les désature davantage. La transition se fait entre $C=0{,}004$
et $C=0{,}012$.

\subsection{La teinte tourne}
Lightroom fait tourner les teintes en même temps qu'il les renforce. Rotation
médiane par famille de teinte, sur 13 scènes :

\begin{center}
\begin{tabular}{lrrr}
\toprule
Famille (teinte Oklab) & $+100$ & $+50$ & $-50$\\
\midrule
rouge (0°) & $+1{,}4^\circ$ & $+0{,}3^\circ$ & $-1{,}5^\circ$\\
orange (30°) & $+3{,}3^\circ$ & $+0{,}9^\circ$ & $-5{,}2^\circ$\\
jaune-orangé (60°) & $+4{,}6^\circ$ & $+2{,}0^\circ$ & $-3{,}6^\circ$\\
vert (150°) & $+0{,}1^\circ$ & $0{,}0^\circ$ & $0{,}0^\circ$\\
cyan (210°) & $-3{,}6^\circ$ & $-2{,}5^\circ$ & $+3{,}9^\circ$\\
bleu-cyan (240°) & $-7{,}7^\circ$ & $-2{,}5^\circ$ & $+5{,}4^\circ$\\
bleu (270°) & $-5{,}1^\circ$ & $-1{,}5^\circ$ & $+1{,}6^\circ$\\
\bottomrule
\end{tabular}
\end{center}

En renforçant, les bleus partent vers le cyan, les oranges vers le jaune ; en
désaturant, le mouvement s'inverse. Le chapitre 5 explique comment le modèle
le reproduit.

\resultat{La vibrance de Lightroom agit sur trois grandeurs à la fois : la
chroma, la clarté et la teinte. Toutes trois dépendent de la saturation, de la
clarté et de la teinte du pixel de départ.}

\chapter{Construire le modèle}

\section{Pourquoi une table plutôt qu'une formule}
Le chapitre 3 a montré une réponse à trois variables, en cloche selon la
clarté, décroissante selon la saturation, avec des particularités par teinte.
On pourrait chercher une formule à quelques constantes. Chaque hypothèse de
forme risque pourtant de reproduire l'erreur des anciens modèles : imposer un
comportement que Lightroom n'a pas.

Une \textbf{table de réponse} évite ce piège : on découpe l'espace des pixels
en cellules et on stocke, pour chacune, la réponse mesurée. C'est la même
démarche que pour les hautes lumières dans \IPR{} : une référence mesurée,
interpolée, plutôt qu'un modèle paramétrique deviné.

\section{Le découpage}
Trois axes, ceux du pixel neutre :
\begin{itemize}
\item saturation ProPhoto $s$ : 12 classes, de bornes
0 ; 0,05 ; 0,1 ; 0,15 ; 0,2 ; 0,3 ; 0,4 ; \ldots ; 0,9 ; 1 (plus fines près du
gris, où la réponse varie vite) ;
\item clarté Oklab $L$ : 9 classes, de bornes 0 ; 0,2 ; 0,3 ; \ldots ; 0,9 ; 1,2 ;
\item teinte Oklab $h$ : 12 secteurs de $30^\circ$.
\end{itemize}
Soit $12\times9\times12=1\,296$ cellules par valeur mesurée et par grandeur
(gain de chroma, gain de clarté, rotation), pour six valeurs du curseur.

\section{Remplir une cellule}
La valeur d'une cellule est la \textbf{médiane}, sur toutes les scènes, des
$\gamma$ (ou $\lambda$, ou $\Delta h$) de ses pixels. La médiane résiste aux
pixels aberrants (bruit, bord d'un objet) mieux que la moyenne.

Une cellule doit contenir au moins 40 pixels pour être jugée fiable. Sinon,
elle reçoit une valeur de repli :
\begin{itemize}
\item pour les gains : la médiane de sa classe (saturation, clarté), toutes
teintes confondues, puis de sa classe de saturation seule ;
\item pour la rotation : la médiane de sa classe (saturation, teinte), puis de
sa teinte seule, puis zéro.
\end{itemize}
La différence est essentielle. Un gain varie peu d'une teinte à l'autre, une
rotation dépend d'abord de la teinte : les bleus tournent dans un sens, les
oranges dans l'autre. Un premier essai qui repliait aussi la rotation « toutes
teintes confondues » donnait à des cellules bleues vides la rotation des
oranges, majoritaires dans le corpus, donc le mauvais sens.

\resultat{Une valeur de repli doit conserver la variable dont la grandeur
dépend le plus.}

\section{Interpoler entre les cellules}
Un pixel tombe rarement au centre d'une cellule. Prendre la valeur de sa
cellule créerait des marches visibles dans un dégradé de ciel. Le modèle
interpole donc entre les centres des cellules voisines, linéairement sur
chaque axe (interpolation \emph{trilinéaire}). Sur un axe, avec deux centres
$x_0<x_1$ et leurs valeurs $v_0,v_1$ :
\begin{equation}
v(x)=(1-t)\,v_0+t\,v_1,\qquad t=\frac{x-x_0}{x_1-x_0}\in[0,1].
\end{equation}
Au-delà du premier ou du dernier centre, la valeur reste constante. La teinte
est un axe circulaire : le secteur centré sur $345^\circ$ a pour voisin celui
centré sur $15^\circ$.

\section{Entre les valeurs mesurées du curseur}
Les tables existent pour $\pm25$, $\pm50$ et $\pm100$. Pour une valeur
intermédiaire $V$, on interpole linéairement entre les deux valeurs mesurées
qui l'encadrent, du même côté de zéro, avec zéro pour origine (réponse nulle).
L'interpolation se fait sur $\gamma$ et $\lambda$, c'est-à-dire en
logarithme : les gains se combinent par multiplication, et cette interpolation
garantit qu'un gain reste positif.

\section{Monotonie : un curseur cohérent}
Une médiane reste une estimation bruitée. Dans quelques cellules peu
peuplées, elle donnait par exemple un gain plus faible à $+50$ qu'à $+25$. Le
photographe verrait alors la couleur \emph{diminuer} en poussant le curseur.

Le modèle impose donc, cellule par cellule, que $\gamma$ croisse avec la
valeur du curseur. Il utilise la \textbf{régression isotone} : la suite
croissante la plus proche, au sens des moindres carrés, des valeurs mesurées.
L'algorithme \emph{Pool Adjacent Violators}~\cite{ayer1955} la calcule en un
passage : chaque fois que deux valeurs voisines violent l'ordre, on les
remplace par leur moyenne, et on recommence si le bloc obtenu viole l'ordre
avec son voisin.

Exemple : les valeurs $(0{,}162\,;\ 0{,}090\,;\ 0{,}924)$ pour $+25$, $+50$,
$+100$ violent l'ordre entre $+25$ et $+50$. On les remplace par leur moyenne :
$(0{,}126\,;\ 0{,}126\,;\ 0{,}924)$.

La contrainte est appliquée séparément de chaque côté de zéro, qui reste fixe.
Elle a modifié 61 cellules sur 1\,296, sans dégrader la validation croisée.

\section{Les quasi-gris}
Les tables sont construites sur des pixels de chroma supérieure à 0,01. Pour
les quasi-gris, on mesure un \textbf{facteur} $f$ qui multiplie le gain prédit
par la table, par classe de chroma et par valeur du curseur. Il est obtenu par
moindres carrés :
\begin{equation}
f=\frac{\sum_i \gamma_i^{\text{LR}}\,\gamma_i^{\text{table}}}{\sum_i \left(\gamma_i^{\text{table}}\right)^2}.
\end{equation}
À $+100$, il vaut 0,74 près du gris et remonte à 1 vers $C=0{,}016$ : le
renforcement est atténué. À $-100$, il vaut 1,38 : la désaturation est
renforcée. À $\pm25$ et $\pm50$, il reste proche de 1.

Un premier essai calculait un seul facteur par sens du curseur. La valeur
$-100$, aux logarithmes très grands, dominait alors l'ajustement et
désaturait trop les quasi-gris à $-25$ et $-50$ : un facteur par valeur a
corrigé ce biais.

Enfin, sous $C=0{,}01$, le gain de clarté et la rotation s'estompent
proportionnellement à $C/0{,}01$. Un gris exact reste exactement gris, et un
presque-gris ne change pas brusquement de teinte, ce qui aurait révélé le
bruit.

\section{La saturation}
La saturation n'utilise pas de table : sa mesure sur 49 scènes suit
Lightroom à moins de 0,05 de variation de chroma. Elle multiplie la chroma
par $1+S/100$ et conserve $L$ et $h$. Combinée à la vibrance, le gain total de
chroma est $(1+S/100)\,e^{\gamma}$.

\section{Exemple déroulé}
Un bleu de ciel, en RGB linéaire sRGB $(0{,}10\,;\ 0{,}20\,;\ 0{,}45)$, à
vibrance $+100$.
\begin{enumerate}
\item Oklab : $L=0{,}583$, $a=-0{,}017$, $b=-0{,}093$, soit $C=0{,}094$ et
$h=259{,}5^\circ$.
\item ProPhoto linéaire : $(0{,}182\,;\ 0{,}197\,;\ 0{,}415)$. La saturation
vaut donc $s=(0{,}415-0{,}182)/0{,}415$, soit $s=0{,}561$.
\item Position dans les tables : $s$ entre les centres 0,55 et 0,65
($t=0{,}11$) ; $L$ entre 0,55 et 0,65 ($t=0{,}33$) ; $h$ entre les secteurs
centrés sur $255^\circ$ et $285^\circ$ ($t=0{,}15$).
\item Lecture interpolée à $+100$ : $\gamma=0{,}541$ (gain $1{,}72$),
$\lambda=-0{,}043$ (clarté $\times0{,}96$), rotation $-4{,}7^\circ$. Le pixel
n'est pas un quasi-gris ($C>0{,}016$) : facteur 1, pas d'estompage.
\item Résultat : $L'=0{,}558$, $C'=0{,}162$ et $h'=254{,}8^\circ$. De retour en RGB
linéaire sRGB, le pixel vaut $(0{,}015\,;\ 0{,}172\,;\ 0{,}628)$.
\end{enumerate}
Le bleu est plus intense, un peu plus sombre, et a tourné vers le cyan.

\chapter{La rotation de teinte}

\section{Un défaut que la mesure de chroma ne voyait pas}
Les premières versions du modèle reproduisaient le gain de chroma de
Lightroom à quelques centièmes près, et pourtant l'œil trouvait la vibrance
d'\IPR{} « moins forte ». Les couleurs ne basculaient pas vers leurs teintes
voisines : le bleu devait tendre vers le cyan. Une mesure qui ne regarde que la
chroma ne peut pas voir ce défaut, puisqu'une rotation conserve la chroma.

\resultat{Une métrique ne valide que ce qu'elle mesure. L'œil a révélé une
dimension, la teinte, que l'indicateur choisi ignorait.}

\section{D'où vient la rotation ?}
Augmenter la saturation dans un espace RGB revient à éloigner le pixel du gris
le long d'une droite de cet espace. Or les droites d'un espace RGB ne sont pas
des lignes de teinte constante pour la vision : en Oklab, leur teinte dérive.
C'est le même phénomène qui fait virer au violet les bleus saturés dans
d'autres espaces perceptuels, ou qui rend les oranges plus jaunes quand on les
pousse.

Les rotations mesurées sont cohérentes avec une vibrance calculée dans un
espace RGB de type ProPhoto : les bleus partent vers le cyan, les oranges vers
le jaune, les verts ne bougent presque pas. Cela reste une interprétation : ce
qui compte pour le modèle, c'est la rotation mesurée, quelle qu'en soit la
cause.

\section{La table de rotation}
La rotation est mesurée et stockée exactement comme les gains : pour chaque
valeur du curseur et chaque cellule (saturation, clarté, teinte), la médiane
de $\Delta h$ sur les scènes. Trois précautions la distinguent :
\begin{itemize}
\item $\Delta h$ est ramené à $]-180^\circ,180^\circ]$ avant toute médiane :
une teinte passant de $359^\circ$ à $1^\circ$ a tourné de $+2^\circ$, pas de
$-358^\circ$ ;
\item le repli des cellules vides conserve la teinte (chapitre 4) ;
\item entre les valeurs du curseur, l'angle est interpolé linéairement, en
degrés, avec zéro à zéro.
\end{itemize}

Dans le module, la rotation s'applique avec la matrice de rotation plane de
l'équation~\eqref{eq:transformation}, après le gain de chroma. Comme gain et
rotation commutent, l'ordre n'a pas d'importance.

\section{Ce que la rotation a corrigé}
Rotation médiane mesurée sur le rendu d'\IPR{} avant et après l'ajout de la
table, comparée à Lightroom, à $+100$ :

\begin{center}
\begin{tabular}{lrrr}
\toprule
Famille & Lightroom & \IPR{} avant & \IPR{} après\\
\midrule
rouge & $+1{,}35^\circ$ & $0^\circ$ & $+1{,}44^\circ$\\
orange & $+3{,}31^\circ$ & $0^\circ$ & $+3{,}30^\circ$\\
jaune-orangé & $+4{,}61^\circ$ & $0^\circ$ & $+4{,}82^\circ$\\
jaune & $+1{,}48^\circ$ & $0^\circ$ & $+1{,}58^\circ$\\
cyan & $-3{,}62^\circ$ & $0^\circ$ & $-4{,}67^\circ$\\
bleu-cyan & $-7{,}67^\circ$ & $0^\circ$ & $-4{,}55^\circ$\\
bleu & $-5{,}13^\circ$ & $0^\circ$ & $-5{,}10^\circ$\\
\bottomrule
\end{tabular}
\end{center}

À $+50$ et $-50$, l'écart de rotation reste inférieur à $0{,}6^\circ$ pour toutes
les familles. Le bleu-cyan à $+100$ reste en retrait ; le chapitre 6 montre
qu'il s'agit surtout d'un grand ciel bleu saturé, sur lequel Lightroom pousse
plus fort que son comportement médian.

\chapter{Valider le modèle}

\section{Deux niveaux de validation}
Un modèle construit sur des mesures les reproduit toujours bien : ce n'est pas
une preuve. Deux contrôles indépendants sont donc menés.
\begin{enumerate}
\item La \textbf{validation croisée} vérifie que le modèle prédit des images
qu'il n'a pas vues. Les 49 scènes sont réparties en deux groupes (une sur
deux) ; on construit les tables sur un groupe et on mesure l'erreur sur
l'autre, puis on inverse.
\item La \textbf{campagne de rendu} mesure le résultat réel d'\IPR{}, de la
lecture du RAW à l'image finale, contre les exports Lightroom. Elle contrôle
tout ce que la validation croisée ignore : conversions, rendu de base,
compression de sortie.
\end{enumerate}

\section{L'indicateur}
Pour chaque scène, on calcule la variation moyenne de chroma
$\overline{C_{\text{variante}}/C_{\text{neutre}}}-1$ sur deux groupes de pixels :
les couleurs peu saturées ($s<0{,}3$) et les couleurs saturées ($s>0{,}5$).
L'écart est la différence entre le modèle (ou \IPR{}) et Lightroom ; on en
retient la médiane sur les scènes, en valeur absolue, et le 90\textsuperscript{e}
centile, qui renseigne sur les scènes difficiles.

\section{Validation croisée}
Écart médian de variation de chroma (peu saturées / saturées) sur les scènes
exclues de la construction :

\begin{center}
\begin{tabular}{lrrr}
\toprule
Valeur & Ancien modèle A1/A2 & Tables & 90\textsuperscript{e} centile des tables\\
\midrule
$-100$ & 0,020 / 0,122 & 0,002 / 0,006 & 0,007 / 0,026\\
$-50$ & 0,047 / 0,030 & 0,003 / 0,005 & 0,007 / 0,020\\
$-25$ & 0,044 / 0,025 & 0,002 / 0,004 & 0,005 / 0,013\\
$+25$ & 0,079 / 0,076 & 0,004 / 0,005 & 0,012 / 0,015\\
$+50$ & 0,088 / 0,127 & 0,012 / 0,012 & 0,030 / 0,035\\
$+100$ & 0,407 / 0,263 & 0,072 / 0,035 & 0,153 / 0,087\\
\bottomrule
\end{tabular}
\end{center}

Les tables divisent l'erreur par 5 à 20 selon la valeur. Le point faible
reste $+100$ sur les couleurs peu saturées, où la réponse de Lightroom est la
plus forte et la plus variable.

\section{Campagne de rendu}
Variation de chroma médiane sur les scènes, Lightroom puis \IPR{}, et écart
médian par scène (campagne \texttt{mesure-couleur-v008}, 637 états mesurés,
aucune exclusion) :

\begin{center}\small
\begin{tabular}{lrrrr}
\toprule
Valeur & Peu saturées & Saturées & Clarté $Y$ & Écart\\
 & LR / \IPR{} & LR / \IPR{} & LR / \IPR{} & peu / saturées\\
\midrule
$-100$ & $-0{,}97$ / $-0{,}97$ & $-0{,}79$ / $-0{,}79$ & 0,88 / 0,87 & 0,004 / 0,011\\
$-50$ & $-0{,}53$ / $-0{,}53$ & $-0{,}40$ / $-0{,}39$ & 0,95 / 0,95 & 0,003 / 0,004\\
$-25$ & $-0{,}28$ / $-0{,}28$ & $-0{,}19$ / $-0{,}19$ & 0,98 / 0,98 & 0,003 / 0,003\\
$+25$ & $+0{,}11$ / $+0{,}11$ & $+0{,}05$ / $+0{,}06$ & 1,00 / 1,00 & 0,014 / 0,003\\
$+50$ & $+0{,}26$ / $+0{,}28$ & $+0{,}11$ / $+0{,}12$ & 1,00 / 1,00 & 0,031 / 0,007\\
$+100$ & $+0{,}90$ / $+0{,}91$ & $+0{,}22$ / $+0{,}24$ & 0,98 / 0,98 & 0,143 / 0,016\\
\bottomrule
\end{tabular}
\end{center}

Avec l'ancien modèle, la même campagne donnait des écarts de 0,31 / 0,09 à
$+100$, 0,14 / 0,06 à $+50$, et 0,01 / 0,13 à $-100$, sans baisse de clarté au
négatif.

\section{Le cas des grands ciels bleus}
Mesuré par famille de teinte, le gain de chroma d'\IPR{} à $+100$ restait plus
faible que celui de Lightroom sur les bleus (1,66 contre 2,78 pour le
bleu-cyan). Trois vérifications ont localisé la cause.

\begin{enumerate}
\item \textbf{Le modèle est bien appliqué.} Sur les pixels bleus d'\IPR{}, la
table prédit un gain de 1,667 et le rendu donne 1,670 : la compression de
sortie ne retire presque rien (elle concerne 25\,\% de ces pixels, et y réduit
le gain de 2,06 à 1,94).
\item \textbf{Le point de départ est identique.} Les bleus neutres d'\IPR{} et de
Lightroom ont la même chroma (0,055 contre 0,056 en médiane), la même
saturation et la même clarté, scène par scène.
\item \textbf{L'écart vient d'une scène.} La chroma finale des bleus à $+100$
est de 0,118 contre 0,124 en médiane (5\,\% de moins), mais la scène 6V1A9141,
un grand ciel bleu saturé qui couvre 82\,\% de l'image, atteint 0,164 contre
0,196 ($-16\,\%$). Elle pèse lourd dans les statistiques par famille.
\end{enumerate}

Sur ce type de ciel, Lightroom pousse plus fort que son comportement médian.
C'est la trace d'une dépendance à l'image que les 2\,\% de variance « scène »
du chapitre 3 laissaient entrevoir. La reproduire demanderait une variable
supplémentaire, par exemple l'étendue de la zone de couleur, ou un traitement
local ; ce n'est pas fait.

\resultat{Un écart global peut cacher un cas particulier très présent dans
l'échantillon. Avant de corriger un modèle, il faut vérifier qu'il est bien
appliqué, que le point de départ est le même, puis isoler les scènes qui
portent l'écart.}

\section{Historique des versions du modèle}
\begin{center}
\begin{tabular}{lp{11cm}}
\toprule
Version & Changement\\
\midrule
v001 & tables brutes de gain de chroma et de clarté\\
v002 & monotonie imposée par régression isotone\\
v003 & facteur quasi-gris, un par sens du curseur\\
v004 & facteur quasi-gris par valeur (corrige $-25$ et $-50$)\\
v005 & table de rotation de teinte\\
v006 & repli de la rotation qui conserve la teinte (version retenue)\\
\bottomrule
\end{tabular}
\end{center}

Chaque version est une preuve JSON horodatée et immuable
(\texttt{modele-vibrance-lr-vNNN}), qui renvoie à celle qu'elle remplace. Les
campagnes de rendu \texttt{mesure-couleur-v004} à \texttt{v008} mesurent
respectivement l'ancien modèle puis les versions v002, v003, v004 et v006.

\chapter{Implémentation, limites et exercices}

\section{Organisation du code}
\begin{longtable}{p{5.6cm}p{9.4cm}}
\toprule Fichier & Rôle\\ \midrule
\small\texttt{scripts/\allowbreak build\_vibrance\_model.py} & Lit les paires Lightroom,
échantillonne, construit les tables, impose la monotonie, mesure le facteur
quasi-gris, valide par validation croisée, écrit l'en-tête C++ et la preuve
JSON. Refuse un état absent ou dupliqué ; ne réécrit jamais une preuve.\\
\small\texttt{src/\allowbreak core/\allowbreak VibranceModel.h} & En-tête généré : axes, tables de gain de
chroma, de gain de clarté et de rotation, facteur quasi-gris. Ne se modifie
pas à la main.\\
\small\texttt{src/\allowbreak core/\allowbreak Vibrance.h/\allowbreak .cpp} & Le module : conversions, lecture
interpolée des tables, transformation du pixel, boucle sur l'image.\\
\small\texttt{src/\allowbreak editor/\allowbreak imageeditor.cpp} & Le cadenceur appelle
\texttt{Vibrance::apply} après les réglages tonals, sans traitement propre.\\
\small\texttt{tests/\allowbreak vibrance\_math\_test.cpp} & Tests unitaires du modèle.\\
\small\texttt{tests/\allowbreak analyse\_couleur.py} & Campagne de mesure du rendu réel
d'\IPR{} contre Lightroom, rendu en mémoire.\\
\bottomrule
\end{longtable}

\section{Le traitement d'une image}
La fonction \texttt{Vibrance::apply} reçoit le master Rec.2020 linéaire en
\texttt{float} sur trois canaux. Elle :
\begin{enumerate}
\item ne fait rien si les deux curseurs sont à zéro, si l'image est vide ou
d'un autre type ;
\item parcourt les lignes en parallèle (\texttt{cv::parallel\_for\_}) : chaque
pixel est indépendant, aucune donnée n'est partagée en écriture ;
\item laisse intact tout pixel gris ($R=G=B$), sans aller-retour matriciel qui
introduirait une erreur d'arrondi ;
\item convertit le pixel en sRGB linéaire, applique la transformation, revient
en Rec.2020.
\end{enumerate}

\section{Ce que vérifient les tests}
\begin{itemize}
\item identités exactes : curseurs à zéro, gris, y compris un gris HDR à 4 ;
\item résultats finis pour des valeurs négatives ou supérieures à 1 ;
\item la chroma ne diminue jamais quand la vibrance augmente ;
\item continuité de part et d'autre de zéro ;
\item saturation $-100$ : gris ; saturation seule : clarté et teinte conservées ;
\item comportement Lightroom : gain $+100$ plausible sur une couleur modérée,
plus faible sur une couleur saturée, réponse progressive, quasi-suppression des
couleurs faibles à $-100$ ;
\item rotation : à $+100$ le bleu tourne vers le cyan, à $-50$ dans l'autre sens ;
\item chemin image : un gris Rec.2020 inchangé, une couleur modifiée, un type
non pris en charge ignoré.
\end{itemize}

\section{Limites}
\begin{itemize}
\item Les tables décrivent Lightroom avec le profil Camera Neutral et la
version de processus courante, sur 49 scènes de trois boîtiers Canon. Un autre
profil ou une autre version de Lightroom demande une nouvelle mesure.
\item Les médianes décrivent le comportement typique. Sur certaines images,
comme un grand ciel bleu saturé, Lightroom s'en écarte ; la part de variance
non expliquée, 5\,\% au positif, existe bien.
\item Les couleurs poussées hors du gamut sRGB sont compressées à la sortie,
comme le reste du rendu.
\item Les quasi-gris et les pixels très sombres sont traités par extrapolation
prudente : les mesures y sont rares et bruitées.
\end{itemize}

\section{Exercices}
\paragraph{Exercice 1 -- Logarithme d'un gain.} Une couleur passe d'une chroma
0,05 à 0,08. Calculer $\gamma$. Quel $\gamma$ ramène ensuite la chroma à 0,05 ?

\paragraph{Exercice 2 -- Rotation.} Un pixel a $a=-0{,}02$ et $b=-0{,}09$.
Calculer sa chroma et sa teinte, puis $(a',b')$ après une rotation de
$-5^\circ$. Vérifier que la chroma est conservée.

\paragraph{Exercice 3 -- Interpolation entre valeurs.} Une cellule donne
$\gamma=0{,}35$ à $+50$ et $\gamma=0{,}91$ à $+100$. Quel gain de chroma à
$+75$ ? À $+20$, sachant que $\gamma=0{,}16$ à $+25$ ?

\paragraph{Exercice 4 -- Régression isotone.} Rendre croissante la suite
$(-3{,}69\,;\ -0{,}82\,;\ -0{,}36)$ pour $-100$, $-50$, $-25$, puis la suite
$(0{,}20\,;\ 0{,}15\,;\ 0{,}10)$ pour $+25$, $+50$, $+100$.

\paragraph{Exercice 5 -- Choix d'une variable.} Pourquoi la chroma absolue
explique-t-elle mal le gain positif, alors que la saturation relative
l'explique bien ? Indice : comparer un bleu sombre et un bleu clair d'aspect
également vif.

\section{Corrigés}
\paragraph{Corrigé 1.} $\gamma=\ln(0{,}08/0{,}05)=\ln1{,}6\approx0{,}47$. Le
retour demande $\gamma=\ln(0{,}05/0{,}08)\approx-0{,}47$ : en logarithme, les
gains inverses sont opposés.

\paragraph{Corrigé 2.} $C=\sqrt{0{,}0004+0{,}0081}\approx0{,}0922$ ;
$h=\operatorname{atan2}(-0{,}09,-0{,}02)\approx-102{,}5^\circ$, soit
$257{,}5^\circ$. Avec $\cos(-5^\circ)\approx0{,}9962$ et
$\sin(-5^\circ)\approx-0{,}0872$ :
$a'=-0{,}02\times0{,}9962-(-0{,}09)(-0{,}0872)\approx-0{,}0278$ et
$b'=-0{,}02\times(-0{,}0872)+(-0{,}09)\times0{,}9962\approx-0{,}0879$.
On trouve $C'=\sqrt{0{,}000773+0{,}00773}\approx0{,}0922$ : la chroma est
conservée, et la teinte devient $252{,}5^\circ$, plus proche du cyan.

\paragraph{Corrigé 3.} À $+75$, $t=0{,}5$ : $\gamma=0{,}63$, gain $e^{0{,}63}\approx1{,}88$.
À $+20$, entre 0 et $+25$, $t=0{,}8$ : $\gamma=0{,}128$, gain $\approx1{,}14$.

\paragraph{Corrigé 4.} La première suite est déjà croissante et négative : elle
est inchangée. La seconde viole l'ordre deux fois. On regroupe les deux
premiers termes en leur moyenne $0{,}175$, qui viole encore l'ordre avec
$0{,}10$ ; on regroupe les trois, de moyenne $0{,}15$. Résultat :
$(0{,}15\,;\ 0{,}15\,;\ 0{,}15)$.

\paragraph{Corrigé 5.} À aspect égal, un bleu sombre a une chroma Oklab plus
faible qu'un bleu clair, parce que la chroma croît avec la luminosité. Classer
par chroma mélange donc des couleurs vives sombres et des couleurs ternes
claires, qui ne réagissent pas pareil. La saturation relative, rapport
indépendant de l'échelle, regroupe les couleurs selon leur « vivacité »,
critère que la vibrance utilise.

\section{Glossaire}
\begin{description}
\item[Chroma $C$] distance au gris dans le plan $(a,b)$ d'Oklab, mesure
absolue de l'intensité.
\item[Saturation $s$] $(\max-\min)/\max$ en ProPhoto linéaire, mesure relative.
\item[Teinte $h$] angle de la couleur dans le plan $(a,b)$.
\item[Gain log $\gamma$] logarithme du rapport des chromas après et avant.
\item[Cellule] case de la table, définie par une plage de saturation, de clarté
et de teinte.
\item[Repli] valeur donnée à une cellule trop peu peuplée.
\item[Régression isotone] suite monotone la plus proche des mesures.
\item[Quasi-gris] pixel de chroma inférieure à 0,016, traité par un facteur.
\item[Validation croisée] mesure de l'erreur sur des données exclues de la
construction du modèle.
\end{description}


\bibliographystyle{plain}
\bibliography{references-vibrance}
\end{document}
